\subsubsection{Pauli Gates}

The \definition{Pauli Gates} are a set of fundamental single-qubit quantum gates that play a central role in quantum computing. They are represented by the three \textbf{Pauli matrices} $X$, $Y$, and $Z$, and can be interpreted as the quantum analogues of basic operations on a qubit.

\highspace
Each Pauli gate corresponds to a \textbf{rotation of $\pi$ radians} around one of the principal axes of the Bloch sphere:
\begin{itemize}
    \item $X$: rotation around the $x$-axis (bit flip).
    \item $Z$: rotation around the $z$-axis (phase flip).
    \item $Y$: rotation around the $y$-axis (bit and phase flip).
\end{itemize}
These operators are \textbf{unitary} ($U^{\dagger} U = U U^{\dagger} = I$), \textbf{Hermitian} ($U = U^{\dagger}$), and \textbf{self-inverse} ($U^2 = I$), and their eigenstates define the three fundamental measurement bases of a qubit. Together, they form the building blocks for describing single-qubit transformations and quantum dynamics. Also, they have eigenvalues $\pm 1$ and satisfy the property:
\begin{equation*}
   X^{2} = Y^{2} = Z^{2} = I
\end{equation*}
\hl{Meaning that applying the same Pauli gate twice returns the qubit to its original state.}

\longline

\paragraph{Rotations Generated by Pauli Operators}

In quantum mechanics, rotations of a qubit are generated by the Pauli operators. A rotation around an axis $\hat{n}$ by an angle $\theta$ is given by:
\begin{equation}\label{eq:rotation-operator-general}
    R_{\hat{n}}(\theta) = e^{-i \frac{\theta}{2} (\hat{n} \cdot \vec{\sigma})}
\end{equation}
Where $\vec{\sigma} = (X, Y, Z)$ are the Pauli matrices. We will see those matrices in more detail in the next sections, but for now we can just note that they are the generators of rotations around the $x$, $y$, and $z$ axes, respectively.

\highspace
In particular:
\begin{align}
    R_x(\theta) &= e^{-i \frac{\theta}{2} X} \\
    R_y(\theta) &= e^{-i \frac{\theta}{2} Y} \\
    R_z(\theta) &= e^{-i \frac{\theta}{2} Z}
\end{align}
Using the identity $\sigma^2 = I$, these can be expanded as:
\begin{equation}\label{eq:rotation-operator-general-expanded}
    R_{\alpha}(\theta)
    =
    \cos\left(\frac{\theta}{2}\right) I
    -
    i \sin\left(\frac{\theta}{2}\right) \sigma_{\alpha}
\end{equation}